%======================================================================
%  Exercise session, Chapter 3 (2 h). Printed by exercise_sheets.tex
%  (questions only) and instructor_manual.tex (answers, instructor notes).
%  No \label here. Codes in ../code/{python,fenics,cast3m}.
%======================================================================
\clearpage
\renewcommand{\sessionchapter}{Chapter 3 -- From continuum to discrete: numerical methods}
\setcounter{spart}{0}
\setcounter{algorithm}{0}\renewcommand{\thealgorithm}{3.\arabic{algorithm}}
\ifanswers
\sessiontitle{Chapter 3 -- The truss in two finite element codes: session sheet with answers}
  {Notes: Chapter~\ref{B-ch:nc}. Codes: FEniCSx (dolfinx, SLEPc), Cast3M; Python for checks.}
\else
\sessiontitle{Chapter 3 -- The truss in two finite element codes: session sheet}
  {Notes: Chapter~\ref{B-ch:nc}. Codes: FEniCSx (dolfinx, SLEPc), Cast3M; Python for checks.}
\fi

\noindent
One structure, three tools, known answers. The two-bar truss of the Chapter~2 session
($EA = L = 1$, half-angle $\varphi$ at the supports) is made of two bar elements with Green
strain. In part~a you assemble its matrices by hand and compute its linear buckling loads.
In part~b you take your first steps in FEniCS: the energy in UFL, linear buckling with SLEPc,
and the nonlinear path with an arc-length step. In part~c you do the same in Cast3M. At the
end of each part, compare with the table below; a code is trusted when it reproduces it.
Bring the notes: Sections~\ref{B-sec:nc-fe}, \ref{B-sec:nc-lba}, \ref{B-sec:nc-riks}
and~\ref{B-sec:nc-detect}.

\begin{center}
\begin{tabular}{@{}lll@{}}
  \toprule
  & $\varphi = 30^\circ$ & $\varphi = 65^\circ$ \\
  \midrule
  first critical point (Chapter~2) & limit point $0.0481$ & bifurcation $0.2434$ \\
  linear buckling (Green strain)   & $0.250$, $0.750$     & $0.324$, $1.489$ \\
  \bottomrule
\end{tabular}
\end{center}

\begin{instructor}
Timing: part~a 30~min, part~b 45~min, part~c 45~min. Install the two codes on the room
machines beforehand (FEniCSx with \texttt{slepc4py}, e.g.\ from conda-forge; Cast3M 2025). Students
who cannot run FEniCS use \texttt{code/python/truss\_fe.py}, which has the same structure.
\untested{The FEniCSx script \texttt{truss\_fenicsx.py} and the Cast3M input \texttt{truss.dgibi}
have not been run (no installation in the build environment): run both once before the
session. The Python codes were run and give the numbers quoted here.}
\end{instructor}

%======================================================================
\spart{The truss as two bar elements, by hand}

\q Write the Green strain $\varepsilon$ and the energy of a bar between nodes $i$ and $j$.
Derive the force on node $j$ and the $2\times2$ block $\mat{k}$ of the element tangent.
\begin{answer}
$\varepsilon = (\vx\cdot\vx - L^2)/(2L^2)$, $\mathcal{E}_e = \tfrac12 EAL\varepsilon^2$, force $N\vx/L$ with
$N = EA\varepsilon$, and $\mat{k} = (EA/L^3)\vx\vx\T + (N/L)\mat{I}$; the element tangent is
$\bigl[\begin{smallmatrix}\mat{k} & -\mat{k}\\ -\mat{k} & \mat{k}\end{smallmatrix}\bigr]$
(Section~\ref{B-sec:nc-fe}).
\end{answer}

\q Assemble $\mat{K}_0$ for the apex. For a unit load, find the linear normal forces and
$\mat{K}_\sigma$. Compute the linear buckling loads and modes at $30^\circ$ and $65^\circ$, and
compare with the first row of the table.
\begin{answer}
$\mat{K}_0 = \diag(2c^2, 2s^2)$, $N = -1/(2s)$ in each bar, $\mat{K}_\sigma = -(1/s)\mat{I}$; loads
$2c^2s$ (sideways) and $2s^3$ (vertical). At $30^\circ$: $0.75$ and $0.25$, against a limit point
at $0.0481$, a factor $5.2$. At $65^\circ$: $0.324$ and $1.489$, against the bifurcation at
$0.2434$, $+33\%$.
\end{answer}

\q Why is linear buckling so poor at $30^\circ$, and what does a negative eigenvalue of
$(\mat{K}_0 + \eta\mat{K}_\sigma)$ mean in general?
\begin{answer}
At $30^\circ$ the truss fails by snap-through after a large prebuckling deflection ($42\%$ of
the rise): the linearisation keeps the initial geometry. A negative $\eta$ is a critical load
for the reversed loading; for a bearing modelled with Haringx's geometric form it is a real
tensile buckling load (Example~\ref{B-ex:nc-bearing}).
\end{answer}

\begin{instructor}
Do not skip the last question: most codes silently discard negative eigenvalues. The bearing
can be shown with Figure~\ref{B-fig:nc-bearing}.
\end{instructor}

%======================================================================
\spart{First steps in FEniCS}

\qc Build the mesh: three points, two interval cells, a vector Lagrange space of degree~1.
Write the bar energy in UFL with the unit vector $\vect{t}_0$ along each element and the
derivative $d\vu/dS = (\operatorname{grad}\vu)\,\vect{t}_0$. Get the residual and the tangent with
\texttt{derivative}. Check $\mat{K}$ at $\vu = \vect{0}$ against part~a.
\begin{answer}
$\varepsilon = \vect{t}_0\cdot d\vu/dS + \tfrac12\abs{d\vu/dS}^2$; at $\vu = \vect{0}$ the apex block of
$\mat{K}$ is $\diag(1.5, 0.5)$ at $30^\circ$. The listing is in Exercise~\ref{B-exo:nc-computer-fenics}.
\end{answer}

\qc Linear buckling with SLEPc: assemble $\mat{K}_0$ (with $1$ on the fixed rows) and
$-\mat{K}_\sigma$ (with $0$ there) and solve the generalised problem with a shift-and-invert
about $0$. Why $0$ on the fixed rows of $-\mat{K}_\sigma$?
\begin{answer}
$0.25$ and $0.75$. The fixed rows then give $\eta = 1/0$, infinitely far from the target, so
they never pollute the spectrum near $0$.
\end{answer}

\qc Follow the nonlinear path with the arc-length step of Algorithm~\ref{alg:nc-arc}, using the
assembled residual and tangent. Locate the limit point from the change of the number of
negative pivots.
\begin{answer}
The path follows $\eta = v(v - s)(v - 2s)$; the pivot count goes $0\to1$ at $0.04811$, mode
$(0,1)$, ${}^1\mat{u}\T\mat{p} \neq 0$: a limit point. The script imports the Python driver: the
algorithm does not care which code assembled $\mat{K}$.
\end{answer}

\begin{algorithm}[ht]
\caption{Arc-length step by bordering (Crisfield's sphere)}\label{alg:nc-arc}
\begin{algorithmic}[1]
\Require converged $(\mat{u}_k,\eta_k)$; previous increment $(\Delta\mat{u}_{\mathrm{old}},\Delta\eta_{\mathrm{old}})$; step $\Delta s$; weight $\psi$
\State solve $\mat{K}(\mat{u}_k)\,\mat{u}_t = \mat{p}$
\State $\Delta\eta \gets \Delta s/\sqrt{\mat{u}_t\T\mat{u}_t + \psi^2\mat{p}\T\mat{p}}$; flip its sign if $\mat{u}_t\T\Delta\mat{u}_{\mathrm{old}} + \psi^2\mat{p}\T\mat{p}\,\Delta\eta_{\mathrm{old}} < 0$
\State $\Delta\mat{u} \gets \Delta\eta\,\mat{u}_t$ \Comment{predictor}
\Repeat \Comment{corrector}
  \State $\mat{r} \gets \mat{r}(\mat{u}_k + \Delta\mat{u}, \eta_k + \Delta\eta)$;\quad $\mat{K} \gets \mat{K}(\mat{u}_k + \Delta\mat{u})$
  \State solve $\mat{K}\,\delta\mat{u}_r = -\mat{r}$ and $\mat{K}\,\delta\mat{u}_p = \mat{p}$ \Comment{one factorisation}
  \State find the roots $\delta\eta$ of $\norm{\Delta\mat{u} + \delta\mat{u}_r + \delta\eta\,\delta\mat{u}_p}^2 + \psi^2(\Delta\eta + \delta\eta)^2\mat{p}\T\mat{p} = \Delta s^2$
  \State keep the root whose new increment is closest in angle to $(\Delta\mat{u},\Delta\eta)$
  \State $\Delta\mat{u} \gets \Delta\mat{u} + \delta\mat{u}_r + \delta\eta\,\delta\mat{u}_p$;\quad $\Delta\eta \gets \Delta\eta + \delta\eta$
\Until{$\norm{\mat{r}} < \epsilon$} \Comment{if no real root or no convergence: halve $\Delta s$}
\State \Return $(\mat{u}_k + \Delta\mat{u}, \eta_k + \Delta\eta)$ and the number of negative pivots of $\mat{K}$
\end{algorithmic}
\end{algorithm}

%======================================================================
\spart{First steps in Cast3M}

\qc Build the truss with two \texttt{barr} elements, \texttt{mate} ($E = 1$) and \texttt{cara}
($\mathrm{SECT} = 1$). With \texttt{opti epsi lineaire}, compute the linear solution for a unit load,
its stress with \texttt{sigm 'LINE'}, and call \texttt{flambage}. Which geometric matrix does
\texttt{ksigma} use for bars?
\begin{answer}
The answer is read off the result: $0.250$ means the Green-strain matrix $(N/L)\mat{I}$;
$0.333$ means $(N/L)(\mat{I} - \vect{e}\vect{e}\T)$, which ignores the axial part (Section~\ref{B-sec:nc-tools}).
Listing: Exercise~\ref{B-exo:nc-computer-castem}.
\end{answer}

\qc Follow the nonlinear path with \texttt{pasapas} and \texttt{'GRANDS\_DEPLACEMENTS'}, imposing the
apex deflection with \texttt{depi}. Recover the load from the reaction and plot it against
$v(v - s)(v - 2s)$. Why does displacement control pass the limit point here?
\begin{answer}
The maximum should be $0.0481$. The deflection $v$ increases monotonically along the whole
path, so each $v$ has exactly one equilibrium: no snap-back. Any difference with the cubic comes
from the strain measure that \texttt{pasapas} uses for bars.
\end{answer}

\qc At $65^\circ$, add a sideways force $10^{-3}$ times the vertical one. What is the maximum
load, and what would the perfect truss do?
\begin{answer}
About $0.2407$ (Chapter~2, and the Python code); the perfect truss bifurcates sideways at
$0.2434$ onto an unstable branch: an unstable symmetric point, two-thirds law.
\end{answer}

\begin{instructor}
Part~c ends the session on the main message: linear buckling, continuation and
imperfections give three different numbers ($0.324$, $0.2434$, $0.2407$ at $65^\circ$), and the
engineer needs to know which question each one answers.
\end{instructor}
